% GENERATED FILE. Edit canonical lessons, then run npm run generate:books.
% canonical-source-hash: fnv1a64:de27a9dd3480ed4c
% canonical-lessons: FR-C165-guerre, FR-C165-paix, FR-C165-armee, FR-C165-soldat, FR-C165-tradition

\chapter{War, Peace, Army, Soldier, Tradition}
\label{ch:fr-war-peace-army-soldier-tradition}
\hlchaptermodality{165}

\begin{chapteropening}
\textbf{By the end of this chapter:} \emph{I can say war, peace, an army, a soldier and a tradition.}

Complete the last lesson of chapter 165: I can say war, peace, an army, a soldier and a tradition.
\end{chapteropening}

\section[la guerre]{la guerre — war}

\label{lesson:FR-C165-guerre}



\begin{quote}
Before the new one: say the French for a butcher's, then the French for a large river.
\end{quote}

\subsection*{You'll want to know: la guerre}
\textbf{la guerre} — \textquotedblleft{}war\textquotedblright{}.

\begin{cousinweb}
\textbf{guerre} $\leftarrow$ Old French \emph{guerre} $\leftarrow$ Medieval Latin \emph{werra} (EtymDB 2.1), from Frankish \emph{*werra}, \textquotedblleft{}strife\textquotedblright{} — borrowed from the Franks; Germanic \emph{w-} became French \emph{gu-}.
\end{cousinweb}

\begin{grammarlens}[title={What you've built}]
One more word for this chapter.
\end{grammarlens}

\subsection*{Your turn}
\begin{itemize}
\raggedright
  \item \emph{Say it:} \emph{la guerre}
  \item \emph{Say it:} \emph{la guerre}, once more
  \item \emph{Say it:} say \emph{la guerre}
  \item \emph{Recall:} say the French for a butcher's, then the French for a large river, then say \emph{la guerre} again
\end{itemize}

\begin{tcolorbox}[breakable,colback=teal!4,colframe=teal!35!black,title={Before you move on}]
What does la guerre mean? (\textquotedblleft{}War\textquotedblright{}.) Say it once more.
\end{tcolorbox}
\section[la paix]{la paix — peace}

\label{lesson:FR-C165-paix}



\begin{quote}
Before the new one: say the French for a large river, then the French for war.
\end{quote}

\subsection*{You'll want to know: la paix}
\textbf{la paix} — \textquotedblleft{}peace\textquotedblright{}.

\begin{cousinweb}
\textbf{paix} $\leftarrow$ Old French \emph{pais} $\leftarrow$ Latin \emph{pax}, \emph{pacis}, \textquotedblleft{}peace\textquotedblright{} (EtymDB 2.1) — inherited: the \emph{pax} behind \emph{payer}.
\end{cousinweb}

\begin{grammarlens}[title={What you've built}]
One more word for this chapter.
\end{grammarlens}

\subsection*{Your turn}
\begin{itemize}
\raggedright
  \item \emph{Say it:} \emph{la paix}
  \item \emph{Say it:} \emph{la paix}, once more
  \item \emph{Say it:} say \emph{la paix}
  \item \emph{Recall:} say the French for a large river, then the French for war, then say \emph{la paix} again
\end{itemize}

\begin{tcolorbox}[breakable,colback=teal!4,colframe=teal!35!black,title={Before you move on}]
What does la paix mean? (\textquotedblleft{}Peace\textquotedblright{}.) Say it once more.
\end{tcolorbox}
\section[l'armée]{l'armée — an army}

\label{lesson:FR-C165-armee}



\begin{quote}
Before the new one: say the French for war, then the French for peace.
\end{quote}

\subsection*{You'll want to know: l'armée}
\textbf{l'armée} — \textquotedblleft{}an army\textquotedblright{}.

\begin{cousinweb}
\textbf{armée} is the past participle of the verb for \textquotedblleft{}to arm\textquotedblright{}, used as a noun: a force that is armed. (EtymDB 2.1 links it to Medieval Latin \emph{armata}, \textquotedblleft{}an armed force\textquotedblright{}.)
\end{cousinweb}

\begin{grammarlens}[title={What you've built}]
One more word for this chapter.
\end{grammarlens}

\subsection*{Your turn}
\begin{itemize}
\raggedright
  \item \emph{Say it:} \emph{l'armée}
  \item \emph{Say it:} \emph{l'armée}, once more
  \item \emph{Say it:} say \emph{l'armée}
  \item \emph{Recall:} say the French for war, then the French for peace, then say \emph{l'armée} again
\end{itemize}

\begin{tcolorbox}[breakable,colback=teal!4,colframe=teal!35!black,title={Before you move on}]
What does l'armée mean? (\textquotedblleft{}An army\textquotedblright{}.) Say it once more.
\end{tcolorbox}
\section[le soldat]{le soldat — a soldier}

\label{lesson:FR-C165-soldat}



\begin{quote}
Before the new one: say the French for peace, then the French for an army.
\end{quote}

\subsection*{You'll want to know: le soldat}
\textbf{le soldat} — \textquotedblleft{}a soldier\textquotedblright{}.

\begin{cousinweb}
\textbf{soldat} $\leftarrow$ Italian \emph{soldato}, \textquotedblleft{}one who is paid\textquotedblright{}, from \emph{soldo}, \textquotedblleft{}pay\textquotedblright{}, $\leftarrow$ Late Latin \emph{solidus}, a gold coin — borrowed. (EtymDB 2.1)
\end{cousinweb}

\begin{grammarlens}[title={What you've built}]
One more word for this chapter.
\end{grammarlens}

\subsection*{Your turn}
\begin{itemize}
\raggedright
  \item \emph{Say it:} \emph{le soldat}
  \item \emph{Say it:} \emph{le soldat}, once more
  \item \emph{Say it:} say \emph{le soldat}
  \item \emph{Recall:} say the French for peace, then the French for an army, then say \emph{le soldat} again
\end{itemize}

\begin{tcolorbox}[breakable,colback=teal!4,colframe=teal!35!black,title={Before you move on}]
What does le soldat mean? (\textquotedblleft{}A soldier\textquotedblright{}.) Say it once more.
\end{tcolorbox}
\section[la tradition]{la tradition — a tradition}

\label{lesson:FR-C165-tradition}



\begin{quote}
Before the new one: say the French for an army, then the French for a soldier.
\end{quote}

\subsection*{You'll want to know: la tradition}
\textbf{la tradition} — \textquotedblleft{}a tradition\textquotedblright{}.

\begin{cousinweb}
\textbf{tradition} $\leftarrow$ Old French \emph{tradicion} $\leftarrow$ Latin \emph{traditio}, \textquotedblleft{}a handing over\textquotedblright{} — a learned borrowing. (EtymDB 2.1)
\end{cousinweb}

\begin{grammarlens}[title={What you've built}]
One more word for this chapter.
\end{grammarlens}

\subsection*{Your turn}
\begin{itemize}
\raggedright
  \item \emph{Say it:} \emph{la tradition}
  \item \emph{Say it:} \emph{la tradition}, once more
  \item \emph{Say it:} say \emph{la tradition}
  \item \emph{Recall:} say the French for an army, then the French for a soldier, then say \emph{la tradition} again
\end{itemize}

\begin{tcolorbox}[breakable,colback=teal!4,colframe=teal!35!black,title={Before you move on}]
What does la tradition mean? (\textquotedblleft{}A tradition\textquotedblright{}.) Say it once more.
\end{tcolorbox}
